% Cuestionario resuelto de la Práctica de Laboratorio 02 — Proyecto OLA.
% Compilar con compilar.sh (dos pasadas, para el total de páginas).
\documentclass[11pt,a4paper]{article}
% Estilo común del informe y del cuestionario de la Práctica de Laboratorio 02 (proyecto OLA).
\usepackage[T1]{fontenc}
\usepackage[utf8]{inputenc}
\usepackage{textcomp}
\DeclareUnicodeCharacter{2212}{\textminus}
\DeclareUnicodeCharacter{2192}{\ensuremath{\rightarrow}}
\DeclareUnicodeCharacter{2264}{\ensuremath{\leq}}
\DeclareUnicodeCharacter{2248}{\ensuremath{\approx}}
\usepackage[spanish,es-tabla,es-nodecimaldot,es-noindentfirst]{babel}
\usepackage[sc]{mathpazo}
\usepackage[scaled=0.92]{helvet}
\renewcommand{\ttdefault}{txtt}
\linespread{1.05}
\usepackage[a4paper,left=2.3cm,right=2.3cm,top=2.7cm,bottom=2.5cm,headsep=0.9cm]{geometry}
\usepackage{microtype}
\usepackage[table]{xcolor}
\usepackage{graphicx}
\usepackage{array,tabularx,longtable,multirow,makecell}
\usepackage{titlesec}
\usepackage{fancyhdr}
\usepackage{lastpage}
\usepackage{enumitem}
\usepackage{listings}
\usepackage[font=small,labelfont={sf,bf},labelsep=colon,skip=6pt]{caption}
\usepackage{pdfpages}
\usepackage{afterpage}
\usepackage[hidelinks,pdfusetitle]{hyperref}
\usepackage{xurl}
\setlength{\emergencystretch}{3em}

% ------------------------------------------------------------------ colores
\definecolor{azul}{HTML}{0F3849}
\definecolor{verdeazul}{HTML}{1F6F80}
\definecolor{fondo}{HTML}{E3EEF0}
\definecolor{fondotenue}{HTML}{F3F7F8}
\definecolor{borde}{HTML}{B9CBCC}
\definecolor{gris}{HTML}{465B62}
\definecolor{rosa}{HTML}{E11D48}
\definecolor{ambar}{HTML}{D97706}
\definecolor{verde}{HTML}{059669}

% ------------------------------------------------------------------ títulos
\titleformat{\section}{\sffamily\bfseries\Large\color{azul}}{\thesection.}{0.6em}{}[{\color{verdeazul}\titlerule[0.6pt]}]
\titleformat{\subsection}{\sffamily\bfseries\large\color{verdeazul}}{\thesubsection}{0.8em}{}
\titleformat{\subsubsection}{\sffamily\bfseries\normalsize\color{azul}}{\thesubsubsection}{0.8em}{}
\titlespacing*{\section}{0pt}{1.6em}{0.9em}
\titlespacing*{\subsection}{0pt}{1.3em}{0.6em}
\setcounter{tocdepth}{2}

% ------------------------------------------------------------------ encabezado y pie
\pagestyle{fancy}
\fancyhf{}
\renewcommand{\headrulewidth}{0.4pt}
\renewcommand{\headrule}{{\color{borde}\hrule width\headwidth height\headrulewidth}}
\fancyhead[L]{\sffamily\footnotesize\color{gris}Calidad de Software · Práctica de Laboratorio 02}
\fancyhead[R]{\sffamily\footnotesize\color{gris}Proyecto OLA}
\fancyfoot[C]{\sffamily\footnotesize\color{gris}Página \thepage\ de \pageref*{LastPage}}
\fancypagestyle{plain}{\fancyhf{}\renewcommand{\headrulewidth}{0pt}%
  \fancyfoot[C]{\sffamily\footnotesize\color{gris}\thepage}}

% ------------------------------------------------------------------ texto
\setlength{\parindent}{0pt}
\setlength{\parskip}{0.55em}
\setlist{itemsep=0.25em,topsep=0.3em}
\newcommand{\ev}[1]{\textsf{\textbf{#1}}}                 % identificador de evidencia
\newcommand{\cod}[1]{\texttt{#1}}                          % código en línea
\newcommand{\negrita}[1]{\textbf{#1}}

% ------------------------------------------------------------------ flotantes
\renewcommand{\topfraction}{0.9}
\renewcommand{\bottomfraction}{0.8}
\renewcommand{\textfraction}{0.07}
\renewcommand{\floatpagefraction}{0.7}
\setcounter{topnumber}{3}\setcounter{bottomnumber}{2}\setcounter{totalnumber}{4}
\makeatletter\setlength{\@fptop}{0pt}\setlength{\@fpbot}{0pt plus 1fil}\setlength{\@fpsep}{16pt}\makeatother
\usepackage[section]{placeins}

% ------------------------------------------------------------------ tablas
\setlength{\arrayrulewidth}{0.4pt}
\arrayrulecolor{borde}
\renewcommand{\arraystretch}{1.25}
\newcolumntype{L}[1]{>{\raggedright\arraybackslash}p{#1}}
\newcolumntype{C}[1]{>{\centering\arraybackslash}p{#1}}
\newcolumntype{Y}{>{\raggedright\arraybackslash}X}
\newcommand{\cab}[1]{\textsf{\bfseries\color{white}#1}}    % celda de cabecera
\newcommand{\filacab}{\rowcolor{azul}}
\newcommand{\filatotal}{\rowcolor{fondo}}
\newcommand{\clave}[1]{\cellcolor{fondo}\textsf{\bfseries\small #1}}

% ------------------------------------------------------------------ fragmentos de código
\lstdefinestyle{ola}{
  language=Python,
  basicstyle=\ttfamily\footnotesize,
  keywordstyle=\color{verdeazul}\bfseries,
  stringstyle=\color{azul},
  commentstyle=\color{gris}\itshape,
  numbers=left, numberstyle=\ttfamily\tiny\color{gris}, numbersep=8pt,
  frame=single, rulecolor=\color{borde}, backgroundcolor=\color{fondotenue},
  xleftmargin=1.6em, framexleftmargin=1.4em, aboveskip=0.8em, belowskip=0.4em,
  columns=fullflexible, keepspaces=true, showstringspaces=false,
  breaklines=true, breakatwhitespace=true, tabsize=4,
  float=htbp, floatplacement=htbp,
  morekeywords={self,None,True,False,Protocol},
}
\lstset{style=ola}
\renewcommand{\lstlistingname}{Fragmento}

% ------------------------------------------------------------------ figuras en página horizontal
% \paginafigura{archivo}{pie}: inserta un PDF A4 horizontal como página propia, con su pie y su número.
\newcommand{\paginafigura}[2]{%
  \includepdf[fitpaper, pagecommand={\thispagestyle{empty}}, picturecommand*={%
    \put(42,13){\parbox{690bp}{\sffamily\footnotesize\color{gris}#2}}%
    \put(800,13){\makebox[0pt][r]{\sffamily\footnotesize\color{gris}Página \thepage\ de \pageref*{LastPage}}}%
  }]{#1}}

% ------------------------------------------------------------------ portada
% \portada{título del documento}{descripción}
\newcommand{\portada}[2]{%
  \hypersetup{pageanchor=false}% la portada no lleva número ni ancla propia
  \begin{titlepage}
    \centering\sffamily\setlength{\parskip}{0pt}
    {\large\bfseries\color{azul}UNIVERSIDAD LA SALLE\par}
    \vspace{0.7cm}
    \includegraphics[width=3.2cm]{figuras/logo_lasalle.png}\par
    \vspace{0.6cm}
    {\large Facultad de Ingenierías y Arquitectura\par}
    \vspace{0.15cm}
    {\large Escuela Profesional de Ingeniería de Software\par}
    \vspace{0.45cm}
    {Curso: Calidad de Software\par}
    \vspace{0.15cm}
    {Docente: Maribel Molina Barriga\par}
    \vspace{1.1cm}
    {\color{verdeazul}\rule{0.62\textwidth}{0.8pt}\par}
    \vspace{0.7cm}
    {\LARGE\bfseries\color{azul}PRÁCTICA DE LABORATORIO 02\par}
    \vspace{0.5cm}
    {\huge\bfseries\color{azul}\begin{minipage}{0.75\textwidth}\centering #1\end{minipage}\par}
    \vspace{0.6cm}
    {\Large\color{verdeazul}\begin{minipage}{0.62\textwidth}\centering Proyecto OLA — Observatorio Litoral de Anomalías térmicas\end{minipage}\par}
    \vspace{0.7cm}
    {\color{verdeazul}\rule{0.62\textwidth}{0.8pt}\par}
    \vspace{0.6cm}
    {\rmfamily\itshape\color{gris}\begin{minipage}{0.66\textwidth}\centering #2\end{minipage}\par}
    \vspace{0.9cm}
    {\bfseries Integrantes\par}
    \vspace{0.35cm}
    {\renewcommand{\baselinestretch}{1.25}\selectfont
    Frederick Mares Graos\par
    Jhordan Huamani Huamani\par
    Jorge Ortiz Castañeda\par
    Miguel Angel Flores Leon\par
    Piero Adrian Delgado Chipana\par}
    \vfill
    Arequipa, octubre de 2026\par
  \end{titlepage}\hypersetup{pageanchor=true}}

\hypersetup{pdftitle={Práctica de Laboratorio 02 — Cuestionario resuelto},
            pdfauthor={Equipo OLA}, pdfsubject={Calidad de Software}}
\titleformat{\section}{\sffamily\bfseries\large\color{azul}}{\thesection.}{0.5em}{}[{\color{verdeazul}\titlerule[0.6pt]}]

\begin{document}

\portada{Cuestionario resuelto}
        {Respuestas al cuestionario de la guía, sustentadas con las mediciones y evidencias del informe técnico de la práctica}

Las respuestas se apoyan en el módulo «Alertas y avisos» del backend de OLA, evaluado en el commit \cod{d5c054f} de la rama \cod{main}, y en las evidencias \ev{E-01} a \ev{E-11} del informe técnico (\cod{CS\_LAB02\_Informe\_OLA.pdf}). Las cifras provienen de las mediciones del 06/10/2026, cuyas salidas en bruto están en la carpeta \cod{evidencias/}.

% =====================================================================
\section{¿Por qué bajo acoplamiento no significa ausencia de relaciones?}

Porque un sistema solo produce valor cuando sus partes colaboran. Bajo acoplamiento significa que las relaciones son pocas, explícitas y dirigidas hacia elementos estables, de modo que un cambio en una parte no obliga a cambiar las demás.

El módulo evaluado ofrece los dos casos (\ev{E-03}, \ev{E-08}). \cod{clock} no tenía ninguna relación ($Ca = 0$, $Ce = 0$) y eso era un defecto (H-04): una abstracción que nadie usa no desacopla nada, y el servicio seguía llamando a \cod{datetime.now(UTC)} en tres lugares. Tras el rediseño, \cod{clock} pasa a $Ca = 2$. En cambio, \cod{notification\_service} depende ahora de más módulos internos que antes (8 en lugar de 5) y aun así está menos acoplado en el sentido que importa: ya no construye sentencias SQL ni lee el reloj del sistema, sino que delega en el repositorio, en el despachador y en la abstracción \cod{Clock}. Lo que se mide no es cuántas flechas hay, sino hacia qué apuntan y cuánto habría que cambiar si el otro extremo cambia.

% =====================================================================
\section{¿Qué atributo mejora con inyección de dependencias?}

Sobre todo la comprobabilidad: si una clase recibe sus colaboradores en lugar de crearlos, en una prueba pueden sustituirse por dobles. También mejora la modificabilidad, porque una implementación puede reemplazarse sin tocar a quien la usa.

En el proyecto, \cod{NotificationDispatcher} recibe por su constructor el almacén de avisos, el registro de canales y el reloj, definidos como protocolos. Gracias a eso sus 12 pruebas unitarias usan un almacén en memoria, un canal que falla a voluntad y un reloj fijo, y se ejecutan en 0.02~s sin PostgreSQL (\ev{E-09}). Antes del cambio, la misma lógica de envío no tenía pruebas unitarias: solo podía comprobarse con las 34 pruebas de integración de \cod{test\_notifications.py}, que tardan entre 12.5 y 12.9~s en cinco ejecuciones porque necesitan la base de datos.

% =====================================================================
\section{¿Cómo se demuestra una mejora en modificabilidad?}

Con un escenario de cambio concreto, implementado sobre los dos diseños con la misma especificación y medido con los mismos indicadores: módulos que hay que modificar, líneas de código existente que se reescriben, complejidad que se añade y pruebas que dejan de pasar. No basta con afirmar que el diseño «es más flexible».

En esta práctica se añadió el mismo canal nuevo, una pasarela por \emph{webhook} siempre activa, sobre ambos diseños (sección 5 del informe, \ev{E-10}). Los módulos existentes modificados bajaron de 6 a 4 y las líneas reescritas, de 13 a 4. El servicio, el repositorio y el router no se tocaron, y la complejidad del bucle de envío se mantuvo en 4, mientras que en el diseño actual subió a 6. La medición también muestra lo que no mejoró: las pruebas que dejan de pasar no disminuyeron (11 en el diseño actual y 12 en el propuesto), porque cuentan avisos y resúmenes que cambian por especificación. La evidencia son los parches de \cod{escenario\_webhook/}, la salida de \mbox{\cod{git diff -{}-numstat}} y la ejecución de la suite sobre cada implementación.

% =====================================================================
\section{¿Cuándo una clase grande puede estar justificada?}

Cuando es grande pero cohesiva: tiene una sola razón para cambiar, baja complejidad, y dividirla solo repartiría el mismo concepto en varios archivos. El tamaño es una señal para revisar, no un defecto en sí mismo.

En OLA, \cod{db/models.py} es el módulo más grande del alcance (229 líneas, 15 clases) y el más referenciado ($Ca = 20$), pero es declarativo: describe tablas y su complejidad máxima es 1 (\ev{E-03}). Otro caso es \cod{Settings}, con 21 campos que representan una única cosa, la configuración de la aplicación. En cambio, \cod{notification\_service}, con 141 líneas, sí debía dividirse: no por su tamaño, sino porque cambiaba por cinco motivos distintos (planificar, persistir, redactar, enviar y marcar la lectura de los avisos; H-01).

% =====================================================================
\section{¿Qué debe contener una evidencia auditable?}

Lo necesario para que otra persona pueda localizarla, entenderla y reproducirla:

\begin{itemize}
  \item un identificador único y la afirmación que respalda;
  \item la versión exacta del artefacto evaluado (aquí, el commit \cod{d5c054f});
  \item la herramienta, su versión y el comando o procedimiento usado;
  \item la fecha y el entorno de ejecución, incluida la configuración que puede afectar el resultado;
  \item el resultado en bruto, o un archivo legible que lo contenga;
  \item quién la obtuvo y sus limitaciones.
\end{itemize}

En el informe cada evidencia, de \ev{E-01} a \ev{E-11}, está numerada, se cita donde se usa y remite a un archivo entregado; las salidas de pytest, ruff, mypy, cobertura, tiempos y métricas están en \cod{evidencias/} con el identificador en su nombre. Las herramientas, sus versiones y la configuración del entorno están en la sección 1.2 del informe; el registro de evidencias, en la 6.2, y las limitaciones, en la 6.3.

\end{document}
